\chapter{Introduction}

\begin{quote}
\textit{``The important thing is not to stop questioning. Curiosity has its own reason for existing.''} \\
\hspace*{\fill}--- Albert Einstein
\end{quote}

Curiosity is widely viewed as a driving force behind learning, both in humans and in artificial systems. In reinforcement learning (RL), this idea is called intrinsic motivation: an internal signal that encourages an agent to seek out new or useful experiences when external feedback is scarce. Yet the question of \textit{how strongly} an agent should listen to its own curiosity remains an open problem. Most existing methods use a single coefficient that must be tuned by hand. This coefficient controls the influence of intrinsic rewards the same way for every state. This thesis takes the view that the importance of curiosity is itself state-dependent. We propose a learned scaling mechanism that adapts the influence of intrinsic motivation to the current situation of the agent.

\section{Background and Research Motivation}
Reinforcement learning has become a main approach for training agents to make decisions in complex environments. This approach has shown strong results in Atari games, Go, and robotic control \cite{mnih2015human,silver2016go,sutton2018rl}. Many of these results, however, depend on dense reward signals that give the agent frequent feedback about good behaviors. In contrast, many real-world environments are sparse-reward. The agent receives useful feedback only after a long sequence of actions. This makes exploration much harder \cite{francois2018deep}.

To handle this problem, many research works add an intrinsic reward on top of the extrinsic reward from the environment \cite{oudeyer2007intrinsic,schmidhuber2010formal}. This intrinsic reward encourages the agent to explore. Count-based methods reward the agent for visiting rarely seen states \cite{bellemare2016count,tang2017exploration,ostrovski2017count}. Curiosity-driven methods build intrinsic rewards from signals such as prediction errors, representation differences, or state-entropy estimates. Examples include the Intrinsic Curiosity Module (ICM) \cite{pathak2017curiosity}, Random Network Distillation (RND) \cite{burda2019rnd}, and Rewarding Impact-Driven Exploration (RIDE) \cite{raileanu2020ride}. These methods improve exploration in many tasks. But they all share one key design choice: how strong the intrinsic reward should be compared to the extrinsic reward.

In common practice, this strength is set by a fixed coefficient. This coefficient must be tuned by hand for each environment. The best value often changes across tasks and even across training stages \cite{chen2022eipo,yuan2023airs}. One global coefficient also cannot tell apart different states. In some states exploration is still useful, but in other states exploration becomes a distraction or even harmful. This problem appears in many intrinsic reward families, from count-based bonuses to ICM and RIDE. This shared weakness leads to a broader research view. Rather than building yet another curiosity module, this thesis studies an adaptive mechanism whose effect can be tested across main intrinsic reward modules.

\section{Research Problem}
This thesis studies the following problem. How does adaptive, state-dependent intrinsic reward scaling affect different intrinsic reward modules? Can state-dependent scaling support useful exploration in sparse-reward reinforcement learning without breaking policy optimization?

More specifically, at each timestep $t$ of an episode, let the agent receive an extrinsic reward $R_t^E$ from the environment and an intrinsic reward $R_t^{I,m}$ produced by intrinsic reward module $m$, where $m$ may denote a count-based bonus, ICM, or RIDE. Standard methods optimize a shaped reward of the form
\begin{equation}
\bar{r}_t = R_t^E + \beta R_t^{I,m},
\end{equation}
where $\beta$ is a fixed hyperparameter. The problem with this form is that the same scaling is applied to every state. But the value of exploration can be very different from one state to another. The main research question of this thesis is therefore:
\begin{quote}
\textit{Can a state-dependent scaling mechanism, learned online from the agent's own experience, improve how representative intrinsic reward modules such as count-based bonuses, ICM, and RIDE contribute to learning in sparse-reward environments?}
\end{quote}

\section{Research Objectives and Scope}
The goal of this thesis is to build and test \textbf{Correlation-Aware Reward Exploration (CARE)}. CARE is a state-dependent scaling mechanism for intrinsic rewards. We also study its effect on main intrinsic reward modules in sparse-reward reinforcement learning. The specific goals are as follows:
\begin{itemize}
    \item analyze why fixed intrinsic reward coefficients are a shared weakness across main intrinsic reward modules;
    \item review count-based exploration, ICM, RIDE, and earlier adaptive reward weighting methods as the basis for state-dependent adaptation;
    \item build CARE as a small scaling layer that works with any intrinsic reward module. CARE learns a state-dependent coefficient $\beta(s_t)$ from a correlation objective that links the weighted intrinsic reward to the extrinsic GAE advantage of the value function;
    \item Implement CARE with Proximal Policy Optimization (PPO) and three main intrinsic reward modules (count-based bonuses, ICM, and RIDE). Each one is compared against a sweep of fixed scalar coefficients;
    \item test the effect of adaptive scaling in sparse-reward MiniGrid environments and discuss what the results mean for different intrinsic reward modules.
\end{itemize}

The scope of this thesis is episodic reinforcement learning tasks with sparse rewards. The study is built around three main intrinsic reward families---count-based bonuses, ICM, and RIDE---combined with PPO as the base policy optimization algorithm. For each module, the adaptive CARE variant is compared against a sweep of fixed scalar coefficients $\beta \in \{0.0005,\, 0.001,\, 0.005,\, 0.01,\, 0.05\}$ and a no-intrinsic PPO baseline.

The experiments use benchmark MiniGrid environments \cite{chevalier2018minigrid} rather than real-world robotics or continuous-control tasks. MiniGrid is a standard testbed for sparse-reward exploration research; the small grid worlds and discrete action space allow fast training with many random seeds, which enables controlled testing of CARE and direct comparison with earlier work. The thesis focuses on experimental evidence and clear methodology rather than a full theoretical proof of convergence or optimality.

\section{Research Methodology}
The work follows a research methodology with four stages. First, the thesis reviews basic RL literature and main intrinsic reward modules to find the common weakness of fixed intrinsic weighting. Second, the thesis proposes CARE as an adaptive scaling mechanism. CARE predicts a positive coefficient $\beta(s_t)$ from the current state. This scaling mechanism can be used with any intrinsic reward module. Third, the method is implemented inside a PPO training loop with three main intrinsic reward modules (count-based bonuses, ICM, and RIDE). Finally, CARE is tested on several sparse-reward tasks in MiniGrid. CARE is compared against PPO without intrinsic rewards and against each module with a swept set of fixed intrinsic coefficients.

The evaluation focuses on sample efficiency, learning stability, robustness across environments, and the behavior of the learned scaling function. This mix of theory background, CARE algorithm design, and experiment results gives an experimental view. The view centers on the effect of CARE on intrinsic reward modules, not on building a new intrinsic reward generator.

\section{Thesis Structure}
The remainder of this thesis is organized as follows. Chapter 2 reviews reinforcement learning in sparse-reward environments, representative intrinsic reward modules, and earlier work on adapting the weight of intrinsic rewards. Additionally, this chapter also identifies the research gap that motivates the proposed method. Chapter 3 presents the proposed methodology. The chapter includes the formal problem statement, the CARE design that works with any module, the advantage-correlation training objective, the three intrinsic module setups (count-based, ICM, RIDE), and the full training algorithm. Chapter 3 also states five formal results (Propositions, a Lemma, and a Corollary) that describe CARE's properties. Chapter 4 describes the experimental setup. This includes the MiniGrid benchmark environments, the implementation details, the evaluation metrics, and the analysis of results from the view of CARE's effect on intrinsic reward modules. The conclusion summarizes the main findings, discusses the limits of the proposed approach, and lists directions for future research. Two appendices complete the thesis. Appendix~A contains the full proofs of the formal results from Chapter 3. Appendix~B lists the network architectures, hyperparameters, and training schedule used in the experiments.
